\section*{अध्याय \ref{ch:g7:what-is-respiration} --- श्वसन क्या है?}

\begin{solution}{exo:g7:what-is-respiration:1}
\omterm{def:g7:what-is-respiration:respiration}{श्वसन} वह गैस-विनिमय है जिससे कोई \omterm{def:g6:exploring-our-environment:components}{सजीव} अपने आस-पास से \omterm{def:g7:what-is-respiration:respiration}{ऑक्सीजन} लेता है और
\omterm{def:g7:what-is-respiration:respiration}{कार्बन डाइऑक्साइड} छोड़ता है। \omterm{def:g7:what-is-respiration:respiration}{श्वासोच्छ्वास} सिर्फ़ छाती की दिखती पंप करती
गतियाँ हैं --- यानी उस विनिमय की सेवा करने का एक तरीक़ा; और विनिमय ख़ुद
हर जीवित \omterm{def:g5:first-look-at-cells:cell}{कोशिका} में चलता है।
\end{solution}

\begin{solution}{exo:g7:what-is-respiration:2}
किसी बंद पात्र में ऑक्सीजन-मापक --- जो दिखाता है कि \omterm{def:g7:what-is-respiration:respiration}{ऑक्सीजन} का स्तर गिर
रहा है; और \omterm{prop:g7:what-is-respiration:shown}{चूने का पानी} --- वह साफ़ तरल जो \omterm{def:g7:what-is-respiration:respiration}{कार्बन डाइऑक्साइड} के गुज़रने
पर दूधिया हो जाता है।
\end{solution}

\begin{solution}{exo:g7:what-is-respiration:3}
यह साबित करने के लिए कि बदलाव जीवित विषय से आया है, न कि पात्र से, दिन
की गरमी से या ख़ुद चूने के पानी से: यानी न्यायी जाँच का जुड़वाँ, जिसमें
अकेला फ़र्क़ जीवन का है।
\end{solution}

\begin{solution}{exo:g7:what-is-respiration:4}
\omterm{prop:g3:balanced-diet:jobs}{ऊर्जा}: हर \omterm{def:g5:first-look-at-cells:cell}{कोशिका} में \omterm{def:g4:journey-of-food:digestion}{पोषक तत्व} \omterm{def:g7:what-is-respiration:respiration}{ऑक्सीजन} के साथ धीरे-धीरे जलाए जाते हैं,
और उनकी \omterm{prop:g3:balanced-diet:jobs}{ऊर्जा} \omterm{def:g5:first-look-at-cells:cell}{कोशिका} के काम के लिए छूटती है। \omterm{def:g7:what-is-respiration:respiration}{कार्बन डाइऑक्साइड} जल चुका
बचा-खुचा है, जिसे ढोकर बाहर कर दिया जाता है।
\end{solution}

\begin{solution}{exo:g7:what-is-respiration:5}
एक चूहा, एक कठ-जूँ, अंकुरित मटर, एक कुकुरमुत्ता, तालाब के घोंघे ---
और जीवित \omterm{def:g6:decomposers-and-soil:soil}{मिट्टी} का एक चम्मच।
\end{solution}

\begin{solution}{exo:g7:what-is-respiration:6}
इसकी जाँच अँधेरे में करो: \omterm{def:g6:exploring-our-environment:conditions}{प्रकाश} से चलने वाला निर्माण रुक जाने पर
मरतबान की \omterm{def:g7:what-is-respiration:respiration}{ऑक्सीजन} गिरती है और \omterm{prop:g7:what-is-respiration:shown}{चूने का पानी} धुँधलाता है --- यानी \omterm{def:g7:what-is-respiration:respiration}{श्वसन},
पर्दा हटाकर।
\end{solution}

\begin{solution}{exo:g7:what-is-respiration:7}
समानताएँ: दोनों \omterm{def:g7:what-is-respiration:respiration}{ऑक्सीजन} खाते हैं और \omterm{def:g7:what-is-respiration:respiration}{कार्बन डाइऑक्साइड} छोड़ते हैं; और
दोनों ईंधन की \omterm{prop:g3:balanced-diet:jobs}{ऊर्जा} छोड़ते हैं। फ़र्क़: \omterm{def:g5:first-look-at-cells:cell}{कोशिका} का जलना धीमा और बिना लौ
का है, हज़ारों छोटे चरणों में; और उसकी \omterm{prop:g3:balanced-diet:jobs}{ऊर्जा} एक ही बार \omterm{def:g6:exploring-our-environment:conditions}{प्रकाश} तथा गरमी
में खोने के बजाय काम के लिए पकड़ ली जाती है।
\end{solution}

\begin{solution}{exo:g7:what-is-respiration:8}
दौड़ता बच्चा (कड़ी पेशीय मेहनत) $>$ पढ़ता बच्चा (शरीर का चुपचाप चलता
ख़र्च) $>$ अंकुरित मटर (तेज़ी से बनता, पर छोटा) $>$ सूखा मटर (सुप्त ---
\omterm{def:g7:what-is-respiration:respiration}{श्वसन} लगभग शून्य पर बेकार)। दर मेहनत के पीछे चलती है।
\end{solution}

\begin{solution}{exo:g7:what-is-respiration:9}
रात भर अंकुरित मटरों के \omterm{def:g7:what-is-respiration:respiration}{श्वसन} ने मरतबान की \omterm{def:g7:what-is-respiration:respiration}{ऑक्सीजन} खा ली और \omterm{def:g7:what-is-respiration:respiration}{कार्बन
डाइऑक्साइड} छोड़ी; और सुबह तक इतनी कम \omterm{def:g7:what-is-respiration:respiration}{ऑक्सीजन} बची कि किसी लौ को जलाए रख
सके। इस बीच \omterm{prop:g7:what-is-respiration:shown}{चूने का पानी} रात भर लगातार धुँधलाता रहा होता।
\end{solution}

\begin{solution}{exo:g7:what-is-respiration:10}
मापकों और चूने के पानी के जालों वाले दो पात्र: एक में ताज़ा जीवित
\omterm{def:g6:decomposers-and-soil:soil}{मिट्टी}, और एक में वही \omterm{def:g6:decomposers-and-soil:soil}{मिट्टी} पकाकर निर्जीव की हुई --- यही अकेला फ़र्क़।
अगर जीवित वाले पात्र की \omterm{def:g7:what-is-respiration:respiration}{ऑक्सीजन} गिरे और उसका \omterm{prop:g7:what-is-respiration:shown}{चूने का पानी} धुँधलाए जबकि
पकाया हुआ जुड़वाँ जितना का तितना रहे, तो यह \omterm{def:g6:decomposers-and-soil:soil}{मिट्टी} के जीवन का ही किया
है।
\end{solution}

\begin{solution}{exo:g7:what-is-respiration:11}
नम अनाज जाग उठता है: \omterm{def:g2:from-seed-to-plant:germination}{अंकुरण} शुरू होता है, और \omterm{def:g7:what-is-respiration:respiration}{श्वसन} फुसफुसाहट से बढ़कर
अंकुरित मटर की रफ़्तार की ओर चला जाता है --- अरबों दाने एक साथ अपने
भंडार जलाते हुए, और खाद के काम करते दलों की तरह ढेर को गरम करते हुए।
गरमाहट यह काम और तेज़ कर देती है: नम ढेर ख़ुद को गरम करके बिगड़ने तक ले
जा सकता है। सूखा अनाज ठंडा सोता है; नम अनाज एक धीमी भट्ठी है।
\end{solution}

\begin{solution}{exo:g7:what-is-respiration:12}
व्यवहार: \omterm{def:g7:what-is-respiration:respiration}{श्वासोच्छ्वास} देखा जा सकता है --- छाती की गतियाँ, मिनट में
पंद्रह बार (\cref{ex:g4:breathing:watch})। रसायन: \omterm{def:g7:what-is-respiration:respiration}{श्वसन} उनमें भी चलता
है जिनकी छाती होती ही नहीं --- मटर, कुकुरमुत्ते, \omterm{def:g6:decomposers-and-soil:soil}{मिट्टी} --- और वह
सिर्फ़ गैस-जाँचों से पकड़ा जाता है। सेवा: \omterm{def:g7:what-is-respiration:respiration}{श्वासोच्छ्वास} इसीलिए है कि
फेफड़ों की हवा नई होती रहे और गैस-अदला-बदली चलती रहे --- और वह ठीक तब
दौड़ता है जब \omterm{def:g5:first-look-at-cells:cell}{कोशिकाओं} का जलना दौड़ता है
(\cref{ex:g4:breathing:running})।
\end{solution}

\begin{solution}{pb:g7:what-is-respiration:1}
\textbf{1.} कोई बदलाव नहीं: \omterm{def:g7:what-is-respiration:respiration}{ऑक्सीजन} जमी हुई, \omterm{prop:g7:what-is-respiration:shown}{चूने का पानी} साफ़। क
निर्जीव \omterm{met:g4:what-plants-need:fairtest}{नियंत्रण} है --- यानी वह निश्चलता जिसके सामने रखकर हर दूसरा
मरतबान पढ़ा जाता है।

\textbf{2.} ख: मुश्किल से नापने लायक़ बदलाव --- सूखे मटर सुप्त हैं और
\omterm{def:g7:what-is-respiration:respiration}{श्वसन} बेकार पड़ा है। ग: साफ़ ऑक्सीजन-गिरावट और धुँधलाहट --- \omterm{def:g2:from-seed-to-plant:germination}{अंकुरण} पूरी
रफ़्तार का निर्माण है, और उस काम का दाम \omterm{def:g7:what-is-respiration:respiration}{श्वसन} में चुकाया जाता है।

\textbf{3.} घ: सबसे तेज़ गिरावट और धुँधलाहट --- गरम, चलता \omterm{prop:g3:sorting-living-things:groups}{स्तनधारी} कड़ा
जलाता है। घंटे इसलिए सीमित हैं कि चूहा अपने मरतबान की \omterm{def:g7:what-is-respiration:respiration}{ऑक्सीजन} खा रहा है;
और प्रयोग हवा कम पड़ने से बहुत पहले ख़त्म होना चाहिए।

\textbf{4.} ङ: \omterm{def:g7:what-is-respiration:respiration}{ऑक्सीजन} जमी हुई या बढ़ती हुई --- उजाले में रखे \omterm{def:g1:plants-around-us:plant}{पौधे} का
निर्माण \omterm{def:g7:what-is-respiration:respiration}{कार्बन डाइऑक्साइड} लेता है और उसके \omterm{def:g7:what-is-respiration:respiration}{श्वसन} के इस्तेमाल से ज़्यादा
\omterm{def:g7:what-is-respiration:respiration}{ऑक्सीजन} छोड़ता है। च: \omterm{def:g7:what-is-respiration:respiration}{ऑक्सीजन} गिरती हुई, \omterm{prop:g7:what-is-respiration:shown}{चूने का पानी} धुँधलाता हुआ ---
अँधेरे में निर्माण रुक जाता है और \omterm{def:g7:what-is-respiration:respiration}{श्वसन}, जो कभी रुका ही नहीं, साफ़ दिख
जाता है।

\textbf{5.} च दिखाता है कि \omterm{def:g1:plants-around-us:plant}{पौधा} हमेशा \omterm{def:g7:what-is-respiration:respiration}{श्वसन} करता है --- \omterm{def:g7:what-is-respiration:respiration}{ऑक्सीजन} भीतर,
\omterm{def:g7:what-is-respiration:respiration}{कार्बन डाइऑक्साइड} बाहर, ठीक चूहे की तरह। और ङ उसी \omterm{def:g7:what-is-respiration:respiration}{श्वसन} को दिन के बड़े
कारोबार के नीचे छिपा लेता है: निर्माण की आमद और उपज उसे दबा देती हैं।
दोनों जुड़वाँ मिलकर इन दो विनिमयों को अलग कर देते हैं।

\textbf{6.} अंकुरित मटर प्रति ग्राम तीन गुना कड़ी मेहनत कर रहा है ---
पूरी रफ़्तार से \omterm{def:g1:plants-around-us:parts}{जड़ें} और \omterm{prop:g2:from-seed-to-plant:cycle}{अंकुर} बना रहा है --- और \omterm{def:g7:what-is-respiration:respiration}{श्वसन} उसी मेहनत के इंजन
की आवाज़ है।

\textbf{7.} चूहे की \omterm{def:g3:bones-and-muscles:muscle}{पेशियाँ} और गरम शरीर पोषक तत्वों को \omterm{def:g7:what-is-respiration:respiration}{ऑक्सीजन} के साथ
जलाते हैं; उस जलने से मरतबान की हवा में \omterm{def:g7:what-is-respiration:respiration}{कार्बन डाइऑक्साइड} छूटती है; वह
फँसी हुई हवा चूने के पानी से गुज़ारी जाती है; और \omterm{def:g7:what-is-respiration:respiration}{कार्बन डाइऑक्साइड} उसे
दूधिया कर देती है।

\textbf{8.} ``हर \omterm{def:g6:exploring-our-environment:components}{सजीव} \omterm{def:g7:what-is-respiration:respiration}{श्वसन} करता है --- \omterm{def:g1:animals-around-us:animal}{जंतु}, \omterm{def:g1:plants-around-us:plant}{पौधा}, जागा हुआ \omterm{def:g2:from-seed-to-plant:seed}{बीज} या
मुश्किल से सोया हुआ --- और \omterm{def:g7:what-is-respiration:respiration}{ऑक्सीजन} लेकर \omterm{def:g7:what-is-respiration:respiration}{कार्बन डाइऑक्साइड} छोड़ता है।''

\textbf{9.} \omterm{def:g7:what-is-respiration:respiration}{ऑक्सीजन} गिरती हुई, \omterm{prop:g7:what-is-respiration:shown}{चूने का पानी} धुँधलाता हुआ --- यानी
कुकुरमुत्ता \omterm{def:g7:what-is-respiration:respiration}{श्वसन} करता है। यह जोड़ता है कि यह नियम \omterm{def:g6:decomposers-and-soil:decomposers}{कवकों} पर भी लागू है:
न \omterm{def:g1:plants-around-us:plant}{पौधा}, न \omterm{def:g1:animals-around-us:animal}{जंतु}, और कोई अपवाद नहीं।

\textbf{10.} क्योंकि गोदाम गोदाम-पैमाने का मरतबान ग ही है: नम रखा अनाज
\omterm{def:g2:from-seed-to-plant:germination}{अंकुरित होने} लगता है और कड़ा \omterm{def:g7:what-is-respiration:respiration}{श्वसन} करता है, और इस तरह भंडार को गरम करके
बिगाड़ देता है --- प्रबंधक के सूखेपन के नियम असल में \omterm{def:g7:what-is-respiration:respiration}{श्वसन} पर क़ाबू
हैं।

\textbf{11.} उसमें मौजूद सब कुछ: तालाब के घोंघे और छोटे \omterm{def:g1:animals-around-us:animal}{जंतु}, कीचड़ के
\omterm{def:g6:decomposers-and-soil:decomposers}{अपघटक}, और अपने अँधेरे घंटों में \omterm{def:g1:plants-around-us:plant}{पौधे} तथा हरी अकेली \omterm{def:g5:first-look-at-cells:cell}{कोशिकाएँ}। यह इसलिए
मायने रखता है कि तालाब के सारे बाशिंदे घुली हुई \omterm{def:g7:what-is-respiration:respiration}{ऑक्सीजन} के एक ही छोटे
भंडार पर टिके हैं --- यही वह अर्थव्यवस्था है जिसे अगले अध्याय खोलते
हैं।

\textbf{12.} सिर्फ़ घ ने --- यानी चूहे ने --- \omterm{def:g7:what-is-respiration:respiration}{श्वासोच्छ्वास} दिखाया:
अपनी गैस-अदला-बदली की सेवा करती दिखती छाती की गतियाँ। और जिस भी मरतबान
में जीवन था उसने \omterm{def:g7:what-is-respiration:respiration}{श्वसन} दिखाया: ख ने हलका, और ग, घ तथा च ने साफ़; ङ ने
दिन के निर्माण के नीचे ढका हुआ, पर \omterm{def:g7:what-is-respiration:respiration}{श्वसन} तब भी चल ही रहा था।
\end{solution}
